\section{Giới thiệu}
\label{sec:introduction}

\subsection{Bối cảnh bài toán}
Trong lịch sử phát triển của ngành Trí tuệ Nhân tạo, năng lực biểu diễn tri thức
hình thức và suy diễn logic chặt chẽ đóng vai trò nền tảng đối với các hệ thống
hỗ trợ ra quyết định đòi hỏi tính đúng đắn cao \cite{russell2020aima}. Nếu
như các bài toán tìm kiếm trên không gian trạng thái cổ điển thường giả định rằng
tác tử được quan sát đầy đủ cấu hình của môi trường ngay từ thời điểm xuất phát,
thì phần lớn các bài toán chẩn đoán và điều tra trong thực tế lại diễn ra trong
điều kiện thông tin không đầy đủ. Trong những môi trường quan sát một phần đó,
tri thức mới chỉ được hé lộ từng bước thông qua chuỗi tương tác tuần tự, đồng
thời mỗi quyết định sai lầm do suy đoán thiếu bằng chứng đều có thể phá vỡ tính
đúng đắn của quá trình lập luận. Chính vì vậy, các trò chơi giải đố suy luận
thuần túy trên lưới hữu hạn đã trở thành môi trường chuẩn mực để kiểm chứng năng
lực của các bộ giải thỏa mãn mệnh đề (SAT) và các tác tử dựa trên cơ sở tri thức
\cite{biere2021handbook}.

Tiêu biểu cho lớp bài toán này là trò chơi suy luận trinh thám
\textbf{Griductive} \cite{griductive2025game}, được phát triển dựa trên cảm hứng
từ trò chơi \textit{Clues by Sam} \cite{cluesbysam2024game}. Bối cảnh trò chơi
đặt ra một bàn cờ vuông kích thước $N \times N$ gồm $N^2$ nghi phạm đứng tại các
tọa độ cột và hàng xác định, trong đó mỗi người mang một thân phận ẩn cố định là
tội phạm hoặc người vô tội. Ở trạng thái khởi tạo ban đầu, chỉ một số ít ô trên
bàn cờ được mở sẵn để cung cấp thân phận thực tế kèm theo lời khai của nhân vật
tương ứng. Quy ước nền tảng tạo nên chiều sâu toán học của trò chơi là mọi lời
khai đã được tiết lộ đều trung thực, bất kể người đưa ra lời khai đó thuộc phe
tội phạm hay phe vô tội. Từ tập dữ kiện ban đầu còn hạn chế, người chơi hoặc tác
tử tự động phải kết hợp các mối quan hệ không gian và định lượng để từng bước lật
mở toàn bộ bàn cờ.

\subsection{Thách thức suy diễn}
Mặc dù luật chơi thoạt nhìn khá trực quan, việc xây dựng một hệ thống giải tự
động hoàn chỉnh lại đặt ra nhiều yêu cầu lý thuyết và kỹ thuật chặt chẽ.
Thách thức cốt lõi nằm ở \textbf{nguyên tắc không suy đoán thiếu căn cứ}
(\textit{no-guessing rule}): tại bất kỳ bước suy diễn $t$ nào với cơ sở tri thức
công khai $KB_t$, một ô chưa mở chỉ được phép kết luận là tội phạm khi mọi mô
hình logic thỏa mãn $KB_t$ đều gán trạng thái tội phạm cho ô đó, và tương tự chỉ
được phép kết luận là vô tội khi không tồn tại mô hình hợp lệ nào cho phép ô đó
là tội phạm. Nếu cả hai giả thuyết trái ngược về một nhân vật đều tương thích với
các dữ kiện hiện hữu, mọi phán quyết đưa ra cho nhân vật ấy đều bị xem là chưa đủ
cơ sở chứng minh và bị hệ thống từ chối, cho dù phán quyết ấy có tình cờ trùng
khớp với đáp án ẩn phía sau.

Bên cạnh yêu cầu chặt chẽ về mặt hệ quả logic, bài toán còn đối mặt với sự bùng
nổ tổ hợp trong quá trình mã hóa tri thức sang dạng chuẩn hội (CNF) cũng như
trong quá trình tìm kiếm nghiệm thỏa mãn \cite{cook1971complexity}. Cụ thể hơn,
các lời khai trong trò chơi không chỉ dừng lại ở những khẳng định đơn lẻ hay quan
hệ đồng nhất giữa hai cá nhân, mà còn bao hàm các ràng buộc đếm số lượng tội phạm
trên từng hàng, từng cột, vùng tám ô lân cận, vùng giao nhau giữa nhiều tập hợp,
hoặc tính chẵn lẻ trên một miền bất kỳ. Khi kích thước miền tham chiếu tăng lên,
phép mã hóa trực tiếp không sử dụng biến phụ trợ rất dễ làm số lượng mệnh đề tăng
nhanh theo hàm tổ hợp hoặc hàm mũ \cite{tseitin1983complexity,sinz2005towards}.
Hơn thế nữa, ở mỗi bước suy diễn trên bàn cờ lớn, tác tử phải liên tục kiểm tra
nhiều ô chưa biết để tìm ra bước đi được suy ra tất yếu, đòi hỏi sự phối hợp
nhịp nhàng giữa cơ chế lan truyền mệnh đề đơn vị hiệu quả và chiến lược tái sử
dụng các mô hình thỏa mãn nhằm cắt giảm các truy vấn dư thừa.

\subsection{Mục tiêu và bố cục}
Xuất phát từ những động lực và thách thức nêu trên, đồ án hướng tới mục tiêu
nghiên cứu, thiết kế và xây dựng hoàn chỉnh một hệ thống giải đố suy diễn mệnh đề
tự động từ tầng biểu diễn toán học đến tầng giao diện trực quan ba chiều. Trước
hết, nhóm tập trung xây dựng kiến trúc điều phối trò chơi có khả năng cô lập chặt
chẽ đáp án gốc khỏi tác tử suy diễn, bảo đảm mọi quyết định và gợi ý đều được dẫn
xuất thuần túy từ dữ kiện công khai. Tiếp nối nền tảng đó, nhóm phát triển bộ mã
hóa CNF tự động cho toàn bộ sáu họ manh mối cốt lõi cùng hai phần mở rộng gồm
phép giao vùng hình học và ràng buộc chẵn lẻ dùng chuỗi phép tuyển loại trừ (XOR)
tuyến tính, kết hợp với bộ giải thỏa mãn mệnh đề DPLL tất định tự cài đặt không
phụ thuộc thư viện ngoài \cite{davis1960computing,davis1962machine}. Song song
với phần lõi thuật toán, nhóm thiết kế không gian bàn điều tra ba chiều với hiệu
ứng đóng dấu phản hồi trực tiếp trên mặt thẻ, đồng thời xây dựng quy trình thực
nghiệm trên các tiến trình độc lập để đánh giá toàn diện hiệu năng hệ thống qua
$14$ màn chơi chuẩn từ cỡ $3 \times 3$ đến cỡ $5 \times 5$.

Phần còn lại của báo cáo được tổ chức theo mạch lập luận gồm sáu chương liên kết
chặt chẽ. Chương \ref{sec:architecture} phân tích đặc trưng môi trường tác tử,
làm rõ kiến trúc phân tầng bảo vệ thông tin ẩn và giới thiệu không gian điều tra
ba chiều. Tiếp theo, Chương \ref{sec:formulation} trình bày phương pháp ánh xạ
biến mệnh đề, đại số vùng không gian và chứng minh từng bước các phép biến đổi
tương đương sang dạng chuẩn hội. Trên cơ sở biểu diễn tri thức đó, Chương
\ref{sec:algorithms} đi sâu vào giải thuật DPLL, định lý suy diễn bằng phản chứng,
chứng minh tính đúng đắn của bộ nhớ đệm mô hình chứng nhân (\textit{witness
model}) và quy trình kiểm định nghiệm duy nhất. Chương \ref{sec:experiments} công
bố kết quả đo đạc thực nghiệm, đánh giá đối chứng thành phần (\textit{ablation
study}), phân tích quy luật tăng trưởng mệnh đề và khảo sát các kịch bản biên
cùng giới hạn kỹ thuật. Cuối cùng, Chương \ref{sec:conclusion} đúc kết các đóng
góp trọng tâm kèm hướng phát triển tiếp theo, trước khi Chương
\ref{sec:contribution} trình bày công tác quản lý đồ án, phân công nhiệm vụ và
tuyên bố minh bạch về việc sử dụng công cụ trí tuệ nhân tạo.
